% Part: proof-theory
% Chapter: sequent-calculus
% Section: admissible-derivable

\documentclass[../../../include/open-logic-section]{subfiles}

\begin{document}

\olfileid[id]{pt}{seq}{adm}

\olsection{Aturan Admisibel dan Aturan yang Dapat Diturunkan}

Banyak hasil dalam teori bukti membahas---atau menggunakan---hasil tentang
apakah segala sesuatu yang dapat kita !!{prove} dengan aturan tertentu juga dapat
kita !!{prove} tanpa aturan tersebut. Teorema eliminasi potongan merupakan contoh
yang paling dikenal. Teorema itu menunjukkan bahwa setiap sekuen yang dapat
kita !!{prove} dengan aturan suatu sistem sekuen beserta~\Cut{} juga dapat
kita !!{prove} tanpa~\Cut. Karena sifat semacam ini sangat penting, sifat itu
mempunyai nama.

\begin{defn}
Suatu aturan~$R$ \emph{admisibel} dalam sistem sekuen jika setiap sekuen yang
dapat dibuktikan dalam sistem tersebut bersama aturan~$R$ juga dapat
dibuktikan tanpa~$R$.
\end{defn}

\begin{ex}\ollabel{ex:land-deriv}
Aturan $\LeftR{\land}$ mempunyai bentuk berbeda dalam \Log{G1c} dan
\Log{G3c}. Dalam \Log{G1c} terdapat dua versi aturan: satu dengan konjung
kiri~$!A$ dari !!{formula} utama~$!A \land !B$ dalam premis, dan satu lagi
dengan konjung kanan~$!B$:
\[
\Axiom$ !A, \Gamma \fCenter \Delta$
\RightLabel{$\LeftR{\land}_1$}
\UnaryInf$ !A \land !B, \Gamma \fCenter \Delta$
\DisplayProof
\qquad
\Axiom$!B, \Gamma \fCenter \Delta$
\RightLabel{$\LeftR{\land}_2$}
\UnaryInf$!A \land !B, \Gamma \fCenter \Delta$
\DisplayProof
\]
Sebaliknya, \Log{G3c} hanya mempunyai satu aturan:
\[
\Axiom$ !A, !B, \Gamma \fCenter \Delta$
\RightLabel{$\LeftR{\land}_3$}
\UnaryInf$ !A \land !B, \Gamma \fCenter \Delta$
\DisplayProof
\]
Sekarang kita dapat bertanya, ``Apakah aturan $\LeftR{\land}_3$ dari
\Log{G3c} admisibel dalam \Log{G1c}?'' Jawabannya ya: setiap inferensi yang
menggunakan $\LeftR{\land}_3$ dapat kita simulasikan secara langsung hanya
dengan aturan \Log{G1c}. Selain $\LeftR{\land}_1$ dan $\LeftR{\land}_2$,
kita juga memerlukan aturan struktural~$\LeftR{\Contraction}$:
\[
\Axiom$!A, !B, \Gamma \fCenter \Delta$
\RightLabel{$\LeftR{\land}_1$}
\UnaryInf$!A \land !B, !B, \Gamma \fCenter \Delta$
\RightLabel{$\LeftR{\land}_2$}
\UnaryInf$!A \land !B, !A \land !B, \Gamma \fCenter \Delta$
\RightLabel{\LeftR{\Contraction}}
\UnaryInf$!A \land !B, \Gamma \fCenter \Delta$
\DisplayProof
\]
\end{ex}

Menyimulasikan inferensi yang menggunakan satu aturan dengan aturan-aturan
lain merupakan salah satu cara menunjukkan bahwa suatu aturan admisibel.
Namun, itu bukan satu-satunya cara, dan beberapa aturan bersifat admisibel
tetapi tidak dapat disimulasikan demikian. Aturan yang dapat disimulasikan
disebut \emph{dapat diturunkan}.

\begin{defn}
Suatu aturan~$R$,
\[
\AxiomC{$S_1$}
\AxiomC{$\dots$}
\AxiomC{$S_n$}
\RightLabel{$R$}
\TrinaryInfC{$S$}
\DisplayProof
\]
disebut \emph{dapat diturunkan} dalam suatu sistem jika terdapat
!!{proof} skematis bagi sekuen~$S$ yang menggunakan aturan dan aksioma sistem
tersebut beserta sekuen-sekuen awal tambahan berbentuk $S_1$, \dots,~$S_n$.
\end{defn}

!!^{proof} yang diberikan dalam \olref{ex:land-deriv} menunjukkan bahwa
$\LeftR{\land}_3$ dapat diturunkan dalam~\Log{G1c}, karena !!{proof}
tersebut merupakan !!{proof} skematis bagi kesimpulan $\LeftR{\land}_3$ dari
premis-premis $\LeftR{\land}_3$ sebagai sekuen awal, dengan hanya memakai
aturan \Log{G1c}. (Bukti itu tidak menggunakan satu pun aksioma \Log{G1c}.)

\begin{prop}\ollabel{prop:lor-G3-adm}
Aturan $\LeftR{\land}$ dan $\RightR{\lor}$ dari~\Log{G3c} dapat diturunkan
dalam~\Log{G1c}.
\end{prop}

\begin{proof}
  Bukti bagi $\LeftR{\land}$ diberikan dalam \olref{ex:land-deriv}. Bukti
  bagi $\RightR{\lor}$ dibiarkan sebagai latihan.
\end{proof}

\begin{prob}
Tunjukkan bahwa aturan $\RightR{\lor}$ dari \Log{G3c} dapat diturunkan
dalam~\Log{G1c}.
\end{prob}

\begin{prob}
Andaikan $\liff$ merupakan !!{operator} terdefinisi: $!A \liff !B$ adalah
singkatan bagi $(!A \lif !B) \land (!B \lif !A)$. Tunjukkan bahwa aturan
\begin{enumerate}
\item \Axiom$\Gamma, !A, !B \fCenter \Delta$
\Axiom$\Gamma \fCenter \Delta, !A, !B$
\RightLabel{\LeftR{\liff}}
\BinaryInf$!A \liff !B, \Gamma \fCenter \Delta$
\DisplayProof
\item
\Axiom$!A, \Gamma \fCenter \Delta, !B$
\Axiom$!B, \Gamma \fCenter \Delta, !A$
\RightLabel{\RightR{\liff}}
\BinaryInf$\Gamma \fCenter \Delta, !A \liff !B$
\DisplayProof
\end{enumerate}
dapat diturunkan dalam (a)~\Log{G1c} dan (b)~\Log{G3c}.
\end{prob}

\begin{ex}\ollabel{ex:weak-adm}
Aturan $\RightR{\Weakening}$ dari \Log{G1c},
\[
\Axiom$ \Gamma \fCenter \Delta$
\RightLabel{\RightR{\Weakening}}
\UnaryInf$ \Gamma \fCenter \Delta, !A$
\DisplayProof
\]
bersifat admisibel dalam~\Log{G3c}. Gagasan buktinya sederhana: andaikan
terdapat !!a{proof}~$\pi$ bagi premis $\Gamma \Sequent \Delta$ dalam
\Log{G3c}. Tambahkan !!{formula}~$!A$ pada sukseden setiap sekuen dalam~$\pi$.
Aksioma tetap merupakan aksioma; inferensi yang benar tetap merupakan
inferensi yang benar. Objek akhirnya bukan !!{formula}, melainkan sekuen
akhir !!{proof} yang baru, yaitu
$\Gamma \Sequent \Delta,!A$.

Namun, aturan $\RightR{\Weakening}$ tidak dapat diturunkan dalam~\Log{G3c}.
Andaikan sebaliknya, yakni andaikan kita dapat menemukan !!a{proof} bagi
$\Gamma \Sequent \Delta,!A$ dengan hanya memakai aksioma dan aturan
\Log{G3c}, serta mungkin sekuen awal tambahan~$\Gamma \Sequent \Delta$.
Jika !!{proof} ini sepenuhnya skematis, sebagaimana disyaratkan agar
$\RightR{\Weakening}$ dapat diturunkan, kita dapat mengubahnya menjadi
!!a{proof} bagi~$\Sequent !A$ dengan menghapus $\Gamma$ dan~$\Delta$. Hal ini
akan mengubah sekuen awal tambahan~$\Gamma \Sequent \Delta$ menjadi sekuen
kosong~$\quad \Sequent \quad$. Karena sekuen kosong sama sekali tidak memuat
!!{formula}, sedangkan semua aturan \Log{G3c} mensyaratkan sekurang-kurangnya
satu !!{formula} aktif dalam premisnya, sekuen kosong tidak dapat menjadi
premis inferensi mana pun dalam !!{proof} skematis ini. Akibatnya, !!{proof}
tersebut hanya dapat bersifat skematis jika tidak benar-benar menggunakan
sekuen tambahan $\Gamma \Sequent \Delta$ sebagai sekuen awal. Dengan kata
lain, bukti itu akan menjadi derivasi skematis bagi
$\Gamma \Sequent \Delta,!A$ yang hanya memakai aksioma dan aturan
\Log{G3c}; dengan demikian, kita akan menunjukkan bahwa setiap sekuen
berbentuk ini mempunyai !!a{proof} dalam~\Log{G3c}. Namun, hal itu mustahil,
karena \Log{G3c} sahih dan karena itu hanya dapat !!{prove} sekuen yang benar
dalam setiap !!{structure}. Dengan mengambil
$\Gamma=\Delta=\emptyset$ dan $!A\ident\lfalse$, \Log{G3c} harus dapat
!!{prove}, misalnya, sekuen $\quad \Sequent \lfalse$, padahal tidak.
\end{ex}

Sekarang kita berikan bukti induktif yang baik bahwa pelemahan merupakan
aturan admisibel dalam~\Log{G3c}. Bahkan, hasil yang sedikit lebih kuat
berlaku, sebagaimana ditunjukkan argumen intuitif tadi: dengan menambahkan
$!A$ pada sisi kanan setiap sekuen dalam~$\pi$, kita tidak mengubah struktur
!!{proof}, dan khususnya tidak menambah tingginya. Karena hal ini penting,
sifat tersebut layak mendapat definisi tersendiri.

\begin{defn}
Suatu aturan~$R$,
\[
\AxiomC{$S_1$}
\AxiomC{$\dots$}
\AxiomC{$S_n$}
\RightLabel{$R$}
\TrinaryInfC{$S$}
\DisplayProof
\]
bersifat \emph{admisibel dengan mempertahankan !!{height}}
(atau \emph{!!{hp}-admisibel}) dalam suatu sistem jika, setiap kali
$\Proves[h_i] S_i$ bagi $i=1$, \dots,~$n$, berlaku $\Proves_m S$ dengan
$m=\max(h_1,\dots,h_n)$.
\end{defn}

Perhatikan bahwa agar suatu aturan bersifat admisibel dengan mempertahankan
!!{height}, cukup bahwa setiap kali premis-premis~$S_i$ mempunyai
!!{proof}s~$\pi_i$ dengan !!{height}~$h_i=\pheight{\pi_i}$, terdapat
!!a{proof}~$\pi$ bagi kesimpulan dengan !!{height}~$\pheight{\pi}$ yang
tidak melebihi maksimum dari~$h_i$. Khususnya, jika hanya ada satu premis
$S_1$, cukup bahwa $\pheight{\pi}\leq\pheight{\pi_1}$. Dalam kasus
$\RightR{\Weakening}$ dan $\LeftR{\Weakening}$, kita bahkan mempunyai
$\pheight{\pi}=\pheight{\pi_1}$, tetapi kesamaan ini tidak diperlukan.

\begin{prop}\ollabel{prop:weak-G3c-adm}
Aturan $\RightR{\Weakening}$ dan $\LeftR{\Weakening}$ dari \Log{G1c}
bersifat admisibel dengan mempertahankan !!{height} dalam~\Log{G3c}.
\end{prop}

\begin{proof}
  Kita menunjukkan bahwa jika $\Log{G3c}\Proves\Gamma\Sequent\Delta$
  dengan !!{proof}~$\pi$ yang mempunyai !!{height}~$\pheight{\pi}=n$, maka terdapat
  \Log{G3c}-!!{proof}~$\pi'$ bagi $\Gamma\Sequent\Delta,!A$ dengan
  $\pheight{\pi'}=n$. Kita membuktikannya melalui induksi pada~$n$. Jika
  $n=0$, maka $\pi$ hanya berupa aksioma $\Gamma\Sequent\Delta$ (yakni,
  suatu $!C$ atomik muncul baik dalam $\Gamma$ maupun~$\Delta$), dan
  $\Gamma\Sequent\Delta,!A$ juga merupakan aksioma. Jika
  $n=\pheight{\pi}>0$, !!{proof}~$\pi$ mempunyai inferensi terakhir. Kita
  membedakan kasus berdasarkan aturan yang digunakan. Kita berikan satu
  contoh saja: andaikan inferensi terakhir adalah $\RightR{\land}$. Maka
  $\Delta=\Delta',!B\land !C$, dan sub-!!{proof}s langsung $\pi_1$ dan
  $\pi_2$ dari inferensi terakhir mempunyai sekuen akhir masing-masing
  $\Gamma\Sequent\Delta',!B$ dan $\Gamma\Sequent\Delta',!C$. Dengan kata
  lain, $\pi$ berbentuk:
  \[
  \AxiomC{}
  \RightLabel{$\pi_1$}
  \Deduce$\Gamma \fCenter \Delta', !B$
  \AxiomC{}
  \RightLabel{$\pi_2$}
  \Deduce$\Gamma \fCenter \Delta', !C$
  \RightLabel{\RightR{\land}}
  \BinaryInf$\Gamma \fCenter \Delta', !B \land !C$
  \DisplayProof
  \]
  Kita juga mempunyai
  $n=\max(\pheight{\pi_1},\pheight{\pi_2})+1$. Jadi,
  $\pheight{\pi_1}<n$ dan $\pheight{\pi_2}<n$, sehingga hipotesis induksi
  berlaku: terdapat !!{proof}s $\pi_1'$ dan $\pi_2'$ bagi
  $\Gamma\Sequent\Delta',!B,!A$ dan $\Gamma\Sequent\Delta',!C,!A$.
  Dengan menggunakan aturan $\RightR{\land}$, kita memperoleh
  !!a{proof}~$\pi'$ bagi $\Gamma\Sequent\Delta',!B\land !C,!A$:
  \[
  \AxiomC{}
  \RightLabel{$\pi_1'$}
  \Deduce$\Gamma \fCenter \Delta', !B, !A$
  \AxiomC{}
  \RightLabel{$\pi_2'$}
  \Deduce$\Gamma \fCenter \Delta', !C, !A$
  \RightLabel{\RightR{\land}}
  \BinaryInf$\Gamma \fCenter \Delta', !B \land !C, !A$
  \DisplayProof
  \]
  Kita mempunyai
  $\pheight{\pi'}=\max(\pheight{\pi_1'},\pheight{\pi_2'})+1
  =\max(\pheight{\pi_1},\pheight{\pi_2})+1=\pheight{\pi}$.
  Kasus bagi $\LeftR{\Weakening}$ bersifat simetris.
\end{proof}

\olref{prop:weak-G3c-adm} sangat berguna. Kita memperkenalkan notasi bagi
penerapan berulangnya: jika $\pi$ merupakan !!a{proof} bagi
$\Gamma\Sequent\Delta$, maka $\pi[\Pi\Sequent\Lambda]$ adalah hasil
penerapan admisibilitas pelemahan di sebelah kiri dengan semua !!{formula}s
dalam~$\Pi$ dan admisibilitas pelemahan di sebelah kanan dengan semua
!!{formula}s dalam~$\Lambda$. Hasilnya adalah !!a{proof} bagi
$\Gamma,\Pi\Sequent\Delta,\Lambda$.

\end{document}
